% GENERATED FILE. Edit canonical lessons, then run npm run generate:books.
% canonical-source-hash: fnv1a64:b4a4faf458a21b12
% canonical-lessons: ES-C374-lapiz, ES-C374-cepillo, ES-C374-guante, ES-C374-vela, ES-C374-nudo

\chapter{Five Things You Hold}
\label{ch:el-lapiz-y-el-cepillo}
\hlchaptermodality{374}

\begin{chapteropening}
\textbf{By the end of this chapter:} \emph{I can say el lápiz, el cepillo, el guante, la vela and el nudo, and name five things that fit in one hand.}

Complete ES-C374-nudo: name five things that fit in one hand, and say one thing about each.
\end{chapteropening}

\section[el lápiz]{el lápiz — a stone you write with}

\label{lesson:ES-C374-lapiz}



\begin{quote}
Say \emph{the skirt}, then \emph{the comb}. (\emph{La falda}; \emph{El peine}.)
\end{quote}

\subsection*{What to know first}
\begin{itemize}
\raggedright
  \item el papel — what it goes onto.
  \item azul — a word from the other half of the same stone.
\end{itemize}

\begin{sounds}
\begin{itemize}
\raggedright
  \item \emph{LÁ-piz}, two beats with the weight on the first. The written accent is there because a word ending in \emph{-z} would otherwise be stressed on its last beat. $\to$ pronunciation reference
\end{itemize}
\end{sounds}

\begin{cousinweb}
\textbf{el lápiz} — \textquotedblleft{}pencil.\textquotedblright{}

Latin \textbf{lapis} meant a stone. That is the whole word: no suffix, no diminutive, nothing meaning \textquotedblleft{}little\textquotedblright{}. A pencil is named for the mineral running down its middle.

Where Spanish picked it up is less settled. The Academy's dictionary takes it straight from Latin. Corominas puts it three centuries later and routes it through Italian, where \emph{lapis} was already a painters' word for a drawing stone.

Then there is a near miss worth slowing down for.

\emph{Lapis lazuli} is the deep blue stone, and its name has two halves that come from opposite ends of the earth:

\begin{itemize}
\raggedright
  \item \textbf{lapis} is the Latin you have right here;
  \item \textbf{lazuli} is Arabic \textbf{lāzaward}, which reached Arabic from Persian and Persian from Sanskrit.
\end{itemize}

Spanish \textbf{azul} descends from the second half. \emph{Lápiz} descends from the first. The two words meet inside one stone's name and share no ancestor at all — a real resemblance, and no relationship.
\end{cousinweb}

\begin{grammarlens}[title={What you've built}]
You can now name the thing you write with: \emph{un lápiz}.
\end{grammarlens}

\subsection*{Your turn}
Say these aloud:

\begin{itemize}
\raggedright
  \item \emph{el lápiz}
  \item \emph{un lápiz}
  \item \emph{un lápiz azul}
\end{itemize}

\begin{tcolorbox}[breakable,colback=teal!4,colframe=teal!35!black,title={Before you move on}]
What did Latin \emph{lapis} mean on its own? (A stone — with no \textquotedblleft{}little\textquotedblright{} in it.) And which half of \emph{lapis lazuli} gave Spanish \emph{azul}? (The second half, the Arabic one.)
\end{tcolorbox}
\section[el cepillo]{el cepillo — named after its back}

\label{lesson:ES-C374-cepillo}



\begin{quote}
Say \emph{the comb}, then \emph{the pencil}. (\emph{El peine}; \emph{El lápiz}.)
\end{quote}

\subsection*{What to know first}
\begin{itemize}
\raggedright
  \item el peine — the one you met a moment ago.
  \item el jabón — what it works with.
\end{itemize}

\begin{sounds}
\begin{itemize}
\raggedright
  \item \emph{ce-PI-llo}, three beats with the weight on the middle one. The \emph{ll} is a \emph{y} sound across nearly all of the Spanish-speaking world. $\to$ pronunciation reference
\end{itemize}
\end{sounds}

\begin{cousinweb}
\textbf{el cepillo} — \textquotedblleft{}brush.\textquotedblright{}

\emph{Cepillo} is a small \textbf{cepo}, and \emph{cepo} is Latin \textbf{cippus}: a post, a stake, a block of wood. The Academy's dictionary puts the diminutive right in the entry.

So the brush is named after its \textbf{back}. The part a Roman would have recognised is the block of wood your hand closes around — not the bristles, which are the part you would probably have named it for.

\emph{Cepo} has never gone away, and it has kept that block in three separate jobs:

\begin{itemize}
\raggedright
  \item the block a trap is built on;
  \item the stocks, the timber a prisoner's ankles went through;
  \item the clamp a city locks onto a badly parked wheel.
\end{itemize}

One piece of wood, three very different centuries, and a brush quietly carrying the oldest of them.
\end{cousinweb}

\begin{grammarlens}[title={What you've built}]
You can now name what you clean with: \emph{un cepillo}.
\end{grammarlens}

\subsection*{Your turn}
Say these aloud:

\begin{itemize}
\raggedright
  \item \emph{el cepillo}
  \item \emph{un cepillo}
  \item \emph{el cepillo y el peine}
\end{itemize}

\begin{tcolorbox}[breakable,colback=teal!4,colframe=teal!35!black,title={Before you move on}]
What did Latin \emph{cippus} mean? (A post or a block of wood.) And which part of a brush does its name point at? (The wooden back, not the bristles.)
\end{tcolorbox}
\section[el guante]{el guante — a Germanic word in Romance clothes}

\label{lesson:ES-C374-guante}



\begin{quote}
Say \emph{the pencil}, then \emph{the brush}. (\emph{El lápiz}; \emph{El cepillo}.)
\end{quote}

\subsection*{What to know first}
\begin{itemize}
\raggedright
  \item la mano — what goes inside it.
  \item el zapato — the other thing you pull onto an end of yourself.
\end{itemize}

\begin{sounds}
\begin{itemize}
\raggedright
  \item \emph{GUAN-te}, two beats with the weight on the first. The \emph{u} is pronounced, so \emph{gua} runs together as one glide. $\to$ pronunciation reference
\end{itemize}
\end{sounds}

\begin{cousinweb}
\textbf{el guante} — \textquotedblleft{}glove.\textquotedblright{}

Nearly everything you have met came down from Latin. This one did not.

The Academy's dictionary traces \textbf{guante} — hedging, with a \emph{perhaps} — through Catalan \textbf{guant} to a Frankish word \textbf{want}, a mitten. Dutch still says \emph{want} for the same thing.

What repays attention is the front of the word. Germanic words beginning in \textbf{w-} do not keep it when they move into Romance. They turn up wearing \textbf{gu-}, and two pairs make the swap audible:

\begin{itemize}
\raggedright
  \item \textbf{guerra} beside English \emph{war};
  \item \textbf{Guillermo} beside English \emph{William}.
\end{itemize}

Say each pair out loud and the \emph{w} and the \emph{gu} land in the same place in your mouth.

One correction while it is fresh. Spanish did not perform that change on \emph{guante}. The word was reshaped on its way in and arrived already wearing the \emph{gu-}. Those pairs show a habit of Romance, not an operation Spanish ran on this particular noun.
\end{cousinweb}

\begin{grammarlens}[title={What you've built}]
You can now name what goes on a hand: \emph{un guante}.
\end{grammarlens}

\subsection*{Your turn}
Say these aloud:

\begin{itemize}
\raggedright
  \item \emph{el guante}
  \item \emph{un guante}
  \item \emph{dos guantes}
\end{itemize}

\begin{tcolorbox}[breakable,colback=teal!4,colframe=teal!35!black,title={Before you move on}]
Is \emph{guante} Latin or Germanic? (Germanic — a Frankish word for a mitten.) And what does a Germanic \emph{w-} become when it moves into Romance? (A \emph{gu-}, as in \emph{guerra} and \emph{Guillermo}.)
\end{tcolorbox}
\section[la vela]{la vela — two words wearing one spelling}

\label{lesson:ES-C374-vela}



\begin{quote}
Say \emph{the brush}, then \emph{the glove}. (\emph{El cepillo}; \emph{El guante}.)
\end{quote}

\subsection*{What to know first}
\begin{itemize}
\raggedright
  \item la luz — what it gives.
  \item la noche — when you need it.
\end{itemize}

\begin{sounds}
\begin{itemize}
\raggedright
  \item \emph{VE-la}, two beats with the weight on the first. The \emph{v} and a \emph{b} are the same sound in Spanish, so this begins the way \emph{bien} does. $\to$ pronunciation reference
\end{itemize}
\end{sounds}

\begin{cousinweb}
\textbf{la vela} — \textquotedblleft{}candle.\textquotedblright{}

Spanish has \textbf{two} words spelled \emph{vela}, and the Academy's dictionary keeps them in separate entries because they are separate words.

\textbf{The candle} comes from the verb \textbf{velar}, and \emph{velar} is Latin \textbf{vigilāre}, \textquotedblleft{}to keep watch.\textquotedblright{} A \emph{vela} is what burns while somebody sits up. Its family in Spanish is \emph{vigilia}, \emph{vigilancia}, \emph{centinela} — watching words, every one.

\textbf{The sail} is a different noun: Latin \textbf{vela}, the plural of \textbf{velum}, \textquotedblleft{}cloth.\textquotedblright{}

It is tempting to knit these together — a sail is stretched cloth, a candle has a wick of cloth, so surely one word. That story is wrong, and wrong in a way worth noticing: it starts from the sail and reasons toward the candle, when the candle came from an entirely different verb.

A useful habit falls out of this. When two Spanish words look identical, the question to ask is not \textquotedblleft{}how are these connected?\textquotedblright{} but \textquotedblleft{}are these one word or two?\textquotedblright{} Here the answer is two, and the dictionary says so on the page.
\end{cousinweb}

\begin{grammarlens}[title={What you've built}]
You can now name what you light: \emph{una vela}.
\end{grammarlens}

\subsection*{Your turn}
Say these aloud:

\begin{itemize}
\raggedright
  \item \emph{la vela}
  \item \emph{una vela}
  \item \emph{una vela en la mesa}
\end{itemize}

\begin{tcolorbox}[breakable,colback=teal!4,colframe=teal!35!black,title={Before you move on}]
Which Latin verb is behind \emph{vela}, the candle? (\emph{Vigilāre}, to keep watch.) And is \emph{vela} meaning \textquotedblleft{}sail\textquotedblright{} the same word? (No — a separate noun, from \emph{velum}, cloth.)
\end{tcolorbox}
\section[el nudo]{el nudo — a vowel nobody can account for}

\label{lesson:ES-C374-nudo}



\begin{quote}
Say \emph{the glove}, then \emph{the candle}. (\emph{El guante}; \emph{La vela}.)
\end{quote}

\subsection*{What to know first}
\begin{itemize}
\raggedright
  \item la cuerda — what you tie it in.
  \item el hilo — the thinner version.
\end{itemize}

\begin{sounds}
\begin{itemize}
\raggedright
  \item \emph{NU-do}, two beats with the weight on the first. The \emph{d} between vowels is soft, close to the \emph{th} of English \emph{this}. $\to$ pronunciation reference
\end{itemize}
\end{sounds}

\begin{cousinweb}
\textbf{el nudo} — \textquotedblleft{}knot.\textquotedblright{}

Latin \textbf{nodus} was a knot. English took it twice, as \textbf{node} and as \textbf{nodule}, and both still mean a knot of something.

Spanish took it too — and the vowel came out wrong.

\emph{Nodus} should have given \emph{\textbf{nodo}}. Spanish has \emph{nodo}, in fact, as a later bookish borrowing. What ordinary speech produced was \emph{nudo}, with a \textbf{u} that no sound change accounts for, so the Academy's dictionary reconstructs a spoken Latin form to bridge the gap.

The obvious guess is sitting right there: Latin also had \textbf{nudus}, \textquotedblleft{}naked,\textquotedblright{} which is where \emph{desnudo} comes from. Two words, one vowel apart, surely they leaned on each other.

Corominas says no, and says it flatly — there is no reason for such a confusion, and Catalan, facing the identical pair, went to some trouble to keep its knot and its naked distinct rather than letting them merge.

So the honest answer is that the \emph{u} is unexplained, and that the tidiest available explanation has been examined and turned down.
\end{cousinweb}

\begin{grammarlens}[title={What you've built}]
You can now name what you tie: \emph{un nudo}.
\end{grammarlens}

\subsection*{Your turn}
Say these aloud:

\begin{itemize}
\raggedright
  \item \emph{el nudo}
  \item \emph{un nudo}
  \item \emph{un nudo en la cuerda}
\end{itemize}

\begin{tcolorbox}[breakable,colback=teal!4,colframe=teal!35!black,title={Before you move on}]
What form should Latin \emph{nodus} have produced in Spanish? (\emph{Nodo} — which exists, as a later borrowing.) And does \emph{nudo} come from \emph{nudus}, \textquotedblleft{}naked\textquotedblright{}? (No; Corominas rejects that.)
\end{tcolorbox}
